\documentclass[10pt,a4paper]{amsart}
\usepackage[margin=19mm]{geometry}
\usepackage{amsmath,amssymb,mathtools}
\usepackage[scr=rsfs,cal=euler]{mathalfa}
\usepackage{xfrac}
\usepackage[unicode,hidelinks,bookmarks=false]{hyperref}
\newcommand{\ie}{i.e.\ }
\newcommand{\cf}{cf.\ }
\newcommand{\resp}{respectively\ }
\newcommand{\Subsection}{\subsection}
\providecommand{\nobreakdash}{\nobreak\hbox{-}}
\setlength{\emergencystretch}{2em}
\multlinegap0pt
\usepackage{bm}
\usepackage{bbm}
\usepackage{dsfont}
\usepackage{xspace}
\usepackage{cancel}
\usepackage{mathtools}
\usepackage{amsthm}
\makeatletter
\@ifundefined{theorem}{\newtheorem{theorem}{Theorem}}{}
\@ifundefined{proposition}{\newtheorem{proposition}[theorem]{Proposition}}{}
\makeatother
\makeatletter
\expandafter\ifx\csname abstract\endcsname\relax
\renewenvironment{abstract}
{\begin{quote}
\noindent\par{\scshape\abstractname.}}
{\medskip\noindent
\end{quote}
}
\fi
\newcommand{\twisten}{\otimes^t\hspace{-1pt}}
\makeatother
\title[BV structure on the Hochschild cohomology of twisted tensor ]{BV structure on the Hochschild cohomology of twisted tensor products}
\author{Matthew Antrobus}
\date{Source: arXiv:2506.23911v1 (CC BY 4.0)}
\begin{document}
\maketitle

\noindent\emph{Source and licence.} BV structure on the Hochschild cohomology of twisted tensor products. Version arXiv:2506.23911v1 (CC BY 4.0), distributed under the Creative Commons Attribution 4.0 International licence, as recorded in the arXiv licence metadata for this exact version and shown on the abstract page https://arxiv.org/abs/2506.23911v1. Full rights notice: licenses/arxiv-2506.23911-CC-BY-4.0.txt.\par\smallskip

\noindent\textit{Editorial scope.} Counted unit: the Proposition whose source label is \texttt{None} in the article (source0.tex lines 198--198, source label \texttt{None}), delivered together with the article's own complete original proof of it (source0.tex lines 199--230, one \texttt{proof} environment of 32 lines). No mathematics is written, repaired or re-derived: statement and proof are the article's own text line for line. Required context retained uncounted: none. Every definition, display and label of those lines is retained unchanged. Omitted from this package and not redistributed: 1--197, 231--901. Nothing inside a retained excerpt is omitted or altered: every retained body is delivered exactly as the article's lines are written, so no omission receipt of any kind accompanies this package. Citations: the retained excerpts carry no citation command at all, so no bibliography entry of the article is transcribed and no reference is redistributed. Reference adaptation: none was needed and none was made; every reference inside the delivered unit resolves inside it, so no unresolved reference or citation is produced. Printed slips kept literal: no printed word, symbol, hypothesis or number of the counted statement or of its proof was altered. Layout and class: the article's own class is \texttt{article} with packages including blindtext, geometry, inputenc, fancyhdr, amsfonts, bbm, amsthm, amsmath, enumerate, mathtools, tikz-cd, thmtools, graphicx, datetime; the wrapper is the lane's pinned \texttt{amsart} editorial wrapper at 10pt on A4, which restates the theorem-like environments the retained text needs and adds the article's own macro definitions that the retained text needs (\texttt{\textbackslash twisten}), with extra wrapper packages (bm, bbm, dsfont, xspace, cancel, mathtools). The delivered numbering is the unit's own. Assets: no retained span contains a figure environment, a graphics-inclusion command, a PDF-page inclusion, a table or a TikZ picture; no image file is copied into this package, the package declares no asset dependency and no image file is redistributed with it. Licence: version arXiv:2506.23911v1 of this article is distributed under the Creative Commons Attribution 4.0 International licence, as recorded in the arXiv licence metadata for this exact version and shown on the abstract page https://arxiv.org/abs/2506.23911v1. A full rights notice is delivered at licenses/arxiv-2506.23911-CC-BY-4.0.txt. Characters: every retained excerpt is pure ASCII, so the delivered engine, XeLaTeX with its default Latin Modern text font, reports no missing character.
\begin{proposition} $\langle-,- \rangle$ makes $R \otimes^t S$ into a Frobenius Algebra. \end{proposition}

\begin{proof} We first must show that $ \langle-,- \rangle$ is a well defined morphism $R \otimes^t S \otimes R \otimes^t S \rightarrow k$. We note that restricting our morphism to the direct summands of the graded decomposition
$$R^{a} \otimes S^{b} \otimes R^{a'} \otimes S^{b'} \rightarrow k$$
our map is simply the composition
$$R^{a} \otimes S^{b} \otimes R^{a'} \otimes S^{b'} \xrightarrow{\cdot t(a',b)} R^{a} \otimes R^{a'} \otimes S^{b} \otimes S^{b'} \xrightarrow{\langle-,-\rangle_R \otimes\langle -,- \rangle_S} k\otimes k \xrightarrow{\mu} k$$
which is clearly linear. 
We next check the Frobenius Condition. This is the requirement that
$$ \langle (a \otimes b) \times (c\otimes d), e \otimes f \rangle = \langle a \otimes b, (c \otimes d) \times (e \otimes f)\rangle $$
and indeed, enduring some algebra
\begin{align*}&\langle ( a \otimes b ) \times^t (c \otimes d ), e \otimes f\rangle\\
= \;\;&\langle t(c,b)(ac \otimes bd), e\otimes f\rangle \\
=\;\;&t(c,b)t(e,bd)\langle ac, e \rangle_R \langle bd, f \rangle_S\\
=\;\;&t(c,b)t(e,b)t(e,d)\langle a, ce\rangle_R \langle b, df \rangle_S \\
= \;\;&t(ce,b)t(e,d)\langle a, ce \rangle_R \langle b, df \rangle_R \\
=\;\; & t(e,d) \langle a \otimes b, ce \otimes df\rangle \\
=\;\; & \langle a \otimes b, (c \otimes d) \times^t (e \otimes f)\rangle
\end{align*}
we see this is the case. \\[5pt]
We now only need to show our inner product is non-degenerate.
We know that both $R$ and $S$ are finite dimensional (as the Frobenius condition forces $R \cong D(R)$), so we may choose bases $\{ e_1,\cdots, e_n\}$ and $\{f_1, \cdots, f_m\}$ respectively, and suppose that these bases are homogeneous. Further, as each of $\langle -,- \rangle_R$ and $\langle -,- \rangle_S$ are non-degenerate, we can choose (left/right) dualizing\footnote{This choice of phrase is to avoid confusion between dual bases, which live in the dual.} bases, $\{\epsilon_1, \cdots, \epsilon_n\}$ and $\{ \phi_1, \cdots, \phi_m\}$ so that
$$\langle \epsilon_i, e_j \rangle_R = \delta_{i,j}1_k \textrm{ (left) } \hspace{20pt} \langle f_i, \phi_j \rangle_S = \delta_{i,j}1_k \textrm{ (right) }$$
To construct these explicitly, we take what is often called the ``dual'' basis of $\{e_1,...e_n\}$ in $D(A)$; i.e, the morphisms
$$d_i(e_j) = \delta_{i,j} 1_k$$
and consider the inverse image of this basis under the isomorphisms $(x \mapsto \langle x, - \rangle_R)$ and $(y \mapsto \langle - , y \rangle_S)$. Note, we have chosen dualizing bases in different coordinates for $R$ and $S$, to make our proof easier. \\[5pt]
Suppose we have $x \in R \twisten S$ such that
$$\langle  x, a \otimes b \rangle = 0 \textrm{ for all } a\otimes b \in R\otimes^t S.$$
Then we know that $\{ \epsilon_i \otimes f_j \}_{i,j}$ form a basis of $R \twisten S$, so write $x = \sum_{i,j} \lambda_{i,j} \epsilon_i \otimes f_j$.
\begin{align*}0= \langle x, e_{i_0} \otimes \phi_{j_0} \rangle&= \sum_{i,j}\lambda_{i,j}\langle \epsilon_i \otimes f_j, e_{i_0} \otimes \phi_{j_0} \rangle \\
&= \sum \lambda_{i,j} t(f_j, e_{i_0})\langle\epsilon_i, e_{i_0}\rangle_R \langle f_j, \phi_{j_0} \rangle_S \\
&= \lambda_{i_0,j_0}t(f_{j_0}, e_i).
\end{align*}
So then, as $t: A \otimes B \rightarrow k^{\times}$, we know that $t(f_{j_0}, e_i)\not=0$, and thus $\lambda_{i_0,j_0}=0$. From here it is clear that $x=0$.
\end{proof}

\end{document}
